% GENERATED FILE. Edit canonical lessons, then run npm run generate:books.
% canonical-source-hash: fnv1a64:162f5ed321daf31b
% canonical-lessons: TA-C522-murram, TA-C522-tinnai, TA-C522-talvaram, TA-C522-tun, TA-C522-cakkatai

\chapter{An Old Village House}
\label{ch:ta-an-old-village-house}
\hlchaptermodality{522}

\begin{chapteropening}
\textbf{By the end of this chapter:} \emph{I can describe the parts of an old house.}

Complete the last lesson of chapter 522: I can describe the parts of an old house.
\end{chapteropening}

\section[\texorpdfstring{\ta{முற்றம்} (muṟṟam)}{முற்றம் (muṟṟam)}]{\ta{முற்றம்} (muṟṟam) — a courtyard}

\label{lesson:TA-C522-murram}



\begin{quote}
Before the new one: say the Tamil for a latch, then the Tamil for a flowerpot.
\end{quote}

\subsection*{You'll want to know: \ta{முற்றம்}}
\textbf{\ta{முற்றம்}} — \emph{muṟṟam} — \textquotedblleft{}a courtyard\textquotedblright{}.

\begin{quote}
\ta{பாட்டியின்} \ta{வீட்டில்} \ta{நடுவில்} \ta{ஒரு} \ta{பெரிய} \ta{முற்றம்} \ta{இருக்கிறது}.
\end{quote}

\begin{quote}
\emph{pāṭṭiyiṉ vīṭṭil naṭuvil oru periya muṟṟam irukkiṟatu.}
\end{quote}

\begin{quote}
Grandmother's house has a big courtyard in the middle.
\end{quote}

An open yard inside a house, under the sky. As a -\ta{அம்} noun: \ta{முற்றத்தில்}.

\begin{cousinweb}
\textbf{\ta{முற்றம்}} (\emph{muṟṟam}) is native Dravidian: DEDR 5016 gives \ta{முற்றம்} (\emph{muṟṟam}), \textquotedblleft{}the courtyard of a house, an open space\textquotedblright{}.
\end{cousinweb}

\begin{grammarlens}[title={What you've built}]
One more word for describing what you lived and felt.
\end{grammarlens}

\subsection*{Your turn}
\begin{itemize}
\raggedright
  \item \emph{Say it:} \emph{muṟṟam}
  \item \emph{Say it:} \emph{muṟṟam}, once more
  \item \emph{Say it:} \emph{pāṭṭiyiṉ vīṭṭil naṭuvil oru periya muṟṟam irukkiṟatu.} — Grandmother's house has a big courtyard in the middle.
  \item \emph{Recall:} say the Tamil for a latch, then the Tamil for a flowerpot, then say \emph{muṟṟam} again
\end{itemize}

\begin{tcolorbox}[breakable,colback=teal!4,colframe=teal!35!black,title={Before you move on}]
What does \ta{முற்றம்} mean? (\textquotedblleft{}A courtyard\textquotedblright{}.) Say it once more.
\end{tcolorbox}
\section[\texorpdfstring{\ta{திண்ணை} (tiṇṇai)}{திண்ணை (tiṇṇai)}]{\ta{திண்ணை} (tiṇṇai) — a raised veranda seat}

\label{lesson:TA-C522-tinnai}



\begin{quote}
Before the new one: say the Tamil for a flowerpot, then the Tamil for a courtyard.
\end{quote}

\subsection*{You'll want to know: \ta{திண்ணை}}
\textbf{\ta{திண்ணை}} — \emph{tiṇṇai} — \textquotedblleft{}a raised veranda seat\textquotedblright{}.

\begin{quote}
\ta{மாலையில்} \ta{தாத்தா} \ta{திண்ணையில்} \ta{உட்கார்ந்து} \ta{பேசுவார்}.
\end{quote}

\begin{quote}
\emph{mālaiyil tāttā tiṇṇaiyil uṭkārntu pēcuvār.}
\end{quote}

\begin{quote}
In the evenings grandfather sits on the veranda and talks.
\end{quote}

The raised platform at the front of a house, where people sit and talk.

\begin{cousinweb}
\textbf{\ta{திண்ணை}} (\emph{tiṇṇai}) is native Dravidian: DEDR 3227 gives \ta{திண்ணை} (\emph{tiṇṇai}) for the raised platform or veranda of a house.
\end{cousinweb}

\begin{grammarlens}[title={What you've built}]
One more word for describing what you lived and felt.
\end{grammarlens}

\subsection*{Your turn}
\begin{itemize}
\raggedright
  \item \emph{Say it:} \emph{tiṇṇai}
  \item \emph{Say it:} \emph{tiṇṇai}, once more
  \item \emph{Say it:} \emph{mālaiyil tāttā tiṇṇaiyil uṭkārntu pēcuvār.} — In the evenings grandfather sits on the veranda and talks.
  \item \emph{Recall:} say the Tamil for a flowerpot, then the Tamil for a courtyard, then say \emph{tiṇṇai} again
\end{itemize}

\begin{tcolorbox}[breakable,colback=teal!4,colframe=teal!35!black,title={Before you move on}]
What does \ta{திண்ணை} mean? (\textquotedblleft{}A raised veranda seat\textquotedblright{}.) Say it once more.
\end{tcolorbox}
\section[\texorpdfstring{\ta{தாழ்வாரம்} (tāḻvāram)}{தாழ்வாரம் (tāḻvāram)}]{\ta{தாழ்வாரம்} (tāḻvāram) — a veranda; a corridor}

\label{lesson:TA-C522-talvaram}



\begin{quote}
Before the new one: say the Tamil for a courtyard, then the Tamil for a raised veranda seat.
\end{quote}

\subsection*{You'll want to know: \ta{தாழ்வாரம்}}
\textbf{\ta{தாழ்வாரம்}} — \emph{tāḻvāram} — \textquotedblleft{}a veranda; a corridor\textquotedblright{}.

\begin{quote}
\ta{மழை} \ta{பெய்தபோது} \ta{நாங்கள்} \ta{தாழ்வாரத்தில்} \ta{நின்றோம்}.
\end{quote}

\begin{quote}
\emph{maḻai peytapōtu nāṅkaḷ tāḻvārattil niṉṟōm.}
\end{quote}

\begin{quote}
When it rained we stood in the veranda.
\end{quote}

The covered walkway under the low edge of a roof.

\begin{cousinweb}
\textbf{\ta{தாழ்வாரம்}} (\emph{tāḻvāram}) begins with the native \ta{தாழ்} (\emph{tāḻ}), \textquotedblleft{}to be low\textquotedblright{} (DEDR 3178), for the low eaves; the origin of the second part, -\ta{வாரம்}, is not settled.
\end{cousinweb}

\begin{grammarlens}[title={What you've built}]
One more word for describing what you lived and felt.
\end{grammarlens}

\subsection*{Your turn}
\begin{itemize}
\raggedright
  \item \emph{Say it:} \emph{tāḻvāram}
  \item \emph{Say it:} \emph{tāḻvāram}, once more
  \item \emph{Say it:} \emph{maḻai peytapōtu nāṅkaḷ tāḻvārattil niṉṟōm.} — When it rained we stood in the veranda.
  \item \emph{Recall:} say the Tamil for a courtyard, then the Tamil for a raised veranda seat, then say \emph{tāḻvāram} again
\end{itemize}

\begin{tcolorbox}[breakable,colback=teal!4,colframe=teal!35!black,title={Before you move on}]
What does \ta{தாழ்வாரம்} mean? (\textquotedblleft{}A veranda; a corridor\textquotedblright{}.) Say it once more.
\end{tcolorbox}
\section[\texorpdfstring{\ta{தூண்} (tūṇ)}{தூண் (tūṇ)}]{\ta{தூண்} (tūṇ) — a pillar}

\label{lesson:TA-C522-tun}



\begin{quote}
Before the new one: say the Tamil for a raised veranda seat, then the Tamil for a veranda.
\end{quote}

\subsection*{You'll want to know: \ta{தூண்}}
\textbf{\ta{தூண்}} — \emph{tūṇ} — \textquotedblleft{}a pillar\textquotedblright{}.

\begin{quote}
\ta{கோயிலின்} \ta{தூண்களில்} \ta{அழகான} \ta{சிலைகள்} \ta{இருக்கின்றன}.
\end{quote}

\begin{quote}
\emph{kōyiliṉ tūṇkaḷil aḻakāṉa cilaikaḷ irukkiṉṟaṉa.}
\end{quote}

\begin{quote}
There are beautiful statues on the temple pillars.
\end{quote}

Plural: \ta{தூண்கள்}, \emph{tūṇkaḷ}.

The plural is \ta{தூண்கள்}. \ta{கம்பம்} (\emph{kampam}) is a thinner post or pole.

\begin{cousinweb}
\textbf{\ta{தூண்}} (\emph{tūṇ}) stands in DEDR's appendix of words shared with Indo-Aryan (DEDR App. 44), beside Sanskrit \emph{sthūṇā}, \textquotedblleft{}a post, a pillar\textquotedblright{} (CDIAL 13774); which way the word travelled is not settled.
\end{cousinweb}

\begin{grammarlens}[title={What you've built}]
One more word for describing what you lived and felt.
\end{grammarlens}

\subsection*{Your turn}
\begin{itemize}
\raggedright
  \item \emph{Say it:} \emph{tūṇ}
  \item \emph{Say it:} \emph{tūṇ}, once more
  \item \emph{Say it:} \emph{kōyiliṉ tūṇkaḷil aḻakāṉa cilaikaḷ irukkiṉṟaṉa.} — There are beautiful statues on the temple pillars.
  \item \emph{Recall:} say the Tamil for a raised veranda seat, then the Tamil for a veranda, then say \emph{tūṇ} again
\end{itemize}

\begin{tcolorbox}[breakable,colback=teal!4,colframe=teal!35!black,title={Before you move on}]
What does \ta{தூண்} mean? (\textquotedblleft{}A pillar\textquotedblright{}.) Say it once more.
\end{tcolorbox}
\section[\texorpdfstring{\ta{சாக்கடை} (cākkaṭai)}{சாக்கடை (cākkaṭai)}]{\ta{சாக்கடை} (cākkaṭai) — a gutter; a sewer}

\label{lesson:TA-C522-cakkatai}



\begin{quote}
Before the new one: say the Tamil for a veranda, then the Tamil for a pillar.
\end{quote}

\subsection*{You'll want to know: \ta{சாக்கடை}}
\textbf{\ta{சாக்கடை}} — \emph{cākkaṭai} — \textquotedblleft{}a gutter; a sewer\textquotedblright{}.

\begin{quote}
\ta{மழையால்} \ta{தெருவின்} \ta{சாக்கடை} \ta{நிறைந்தது}.
\end{quote}

\begin{quote}
\emph{maḻaiyāl teruviṉ cākkaṭai niṟaintatu.}
\end{quote}

\begin{quote}
Because of the rain the street gutter filled up.
\end{quote}

An open drain along a street. The water in it is \ta{சாக்கடை} \ta{நீர்}.

\begin{cousinweb}
The origin of \textbf{\ta{சாக்கடை}} (\emph{cākkaṭai}) is not settled: no dictionary we can check gives it.
\end{cousinweb}

\begin{grammarlens}[title={What you've built}]
One more word for describing what you lived and felt.
\end{grammarlens}

\subsection*{Your turn}
\begin{itemize}
\raggedright
  \item \emph{Say it:} \emph{cākkaṭai}
  \item \emph{Say it:} \emph{cākkaṭai}, once more
  \item \emph{Say it:} \emph{maḻaiyāl teruviṉ cākkaṭai niṟaintatu.} — Because of the rain the street gutter filled up.
  \item \emph{Recall:} say the Tamil for a veranda, then the Tamil for a pillar, then say \emph{cākkaṭai} again
\end{itemize}

\begin{tcolorbox}[breakable,colback=teal!4,colframe=teal!35!black,title={Before you move on}]
What does \ta{சாக்கடை} mean? (\textquotedblleft{}A gutter; a sewer\textquotedblright{}.) Say it once more.
\end{tcolorbox}
